\documentclass{beamer}
\usetheme[shownavsym]{AAUsimple}
\usefonttheme{structurebold}
%%%%%%%%%%%%%%%%%%%%%%%%%%%%%%%%%%
\usepackage[utf8]{inputenc}
\usepackage[dutch]{babel}
\usepackage[T1]{fontenc}
\usepackage{tikz}
\usetikzlibrary{positioning}
\usepackage{hyperref}
\usepackage{microtype}
\definecolor{darkblue}{RGB}{33,26,82}
\definecolor{lightblue}{RGB}{194,193,204}
\usepackage{array,multirow, booktabs}%%Tabellenmaker 1.Array voor horizontale lijn en arraystretch, 2. Multirow en Multicolumn voor merging rijen en columnen respectievelijk, 3.Booktabs voor code commands in tabel.
%%%%%%%%%%%%%%%%%%%%%%%%%%%%%%%%%%
% Gekleurde hyperlinks
\newcommand{\chref}[2]{%
	\href{#1}{{\usebeamercolor[bg]{AAUsidebar}#2}}%
}
\title[PKU]{PKU - Nomenclatuur mutaties}

\date{\today}

\author[Edon Namani, et. al]
{%
	Edon Namani\\
	\href{mailto:e.namani@student.rug.nl}{{\tt e.namani@student.rug.nl}}
}

\institute[%
	Faculteit Gezondheidswetenschappen\\
	Rijksuniversiteit Groningen\\
	Nederland
]
{%
	Faculteit Gezondheidswetenschappen\\
	Rijksuniversiteit Groningen\\
	Nederland

}

\pgfdeclareimage[height=1.5cm]{titlepagelogo}{AAUgraphics/semper_invicta.pdf}
\titlegraphic{%
	\pgfuseimage{titlepagelogo}
}

\begin{document}
% Titel
{\aauwavesbg%
	\begin{frame}[plain,noframenumbering]
	\titlepage
\end{frame}}
%%%%%%%%%%%%%%%

%TOC
\begin{frame}{Agenda}{}
	\tableofcontents
\end{frame}
%%%%%%%%%%%%%%%
\section{nomenclatuur}
\begin{frame}{Syntax}
	Mutaties kunnen op zowel gen- als proteïnniveau worden weergeven. Het eerstgenoemde heeft betrekking op de nucleotidesequenctie en het laatstgenoemde op aminozuursequentie.
	
	\vspace{\baselineskip}
	\centering
		\begin{tikzpicture}
			\node[fill=darkblue,text=white] {[niveau].[verandering op niveau]};
		\end{tikzpicture}
		\begin{itemize}
		\item niveau
			\begin{itemize}
			\item c $\rightarrow$ gen
			\item p $\rightarrow$ prote\"ine
			\item g $\rightarrow$ genoom
			\end{itemize}
		\item verandering $\rightarrow$ missense, nonsense, frameshift, etc.
		\end{itemize}
	
\end{frame}

\subsection{Nummering van nucleotides}
\begin{frame}{Nummering van nucleotides}
	\begin{tikzpicture}[text=white,inner sep=0, annotation/.style={anchor=####1,fill=darkblue,inner sep=2pt},
		annotation/.default=north]
		\fill[gray] (0,0) rectangle (2,.5);
		\fill[darkblue] (2,0) rectangle (5,.5);
		\fill[lightblue] (5,0) rectangle (7,.5);
		\fill[darkblue] (7,0) rectangle (9,.5);
		\fill[gray] (9,0) rectangle (10.8,.5);
		\node at (1,.25) {5'UTR};
		\node at (3.5,.25) {Exon 1};
		\node at (6,.25) {Intron 1};
		\node at (8,.25) {Exon 2};
		\node at (9.9,.25) {3'UTR};
	\draw[very thick] (2,0) -- +(0,-1) node[annotation] {c.1} (1,0.5) -- +(0,1) node[annotation=south] {c.-10} (5,0) -- +(0,-1) node[annotation] {c.30}
		(5.25,.5) -- + (0,1) node[annotation=south] {c.30+5} (6.75,.5) -- +(0,1) node[annotation=south] {c.31-5} (7,0) -- +(0,-1) node[annotation] {c.31} (9,0) -- +(0,-1) node[annotation] {c.50} (10,.5)  -- +(0,1) node[annotation=south] {c*10};
\end{tikzpicture}
\end{frame}

\subsection{syntax van mutatie}
\begin{frame}{Nucleotide $\rightarrow$ Prote\"ine}
	\begin{tabular}{*3l}
		\toprule
		Mutatie & Nucleotide verandering & Prote\"ine verandering\\
		\midrule
		missense  & c.400A>C & p.Ile134Arg\\
		nonsense  & c.101A>T & p.Phe134*\\
		frameshift& c.400\_403del & p.Arg134Tyrfs*5\\
		in-frame  & c.200\_202dup & p.Ala67dup\\
		silent    & c.165A>G & p.Phe55=\\
		splicing  & c.170+10T>G & p.?\\
		\bottomrule
	\end{tabular}
\end{frame}

{\aauwavesbg
\begin{frame}[plain,noframenumbering]
	\finalpage{\href{https://docs.google.com/presentation/d/1wWC_u3hx8rsNSXKGCjr6AD3-wVgjlJHiZjDQSCHe08o/edit\#slide=id.g421dd203c7_0_5}{Terug naar Google Slides}
		}
\end{frame}
}
\end{document}
